% 人工对照底本逐节点转写；独立TikZ重绘。

\begin{tikzpicture}[x=1cm,y=-1cm,every node/.style={anchor=west,align=left,inner sep=2pt,font=\fontsize{10}{13.0}\selectfont},edge/.style={draw=black!65,line width=.45pt}]
\node[text width=2.500cm,draw=black!65] (n1) at (0.000,1.092) {第一步：处分行为};
\node[text width=0.500cm,draw=black!65] (n2) at (2.920,0.287) {无};
\node[text width=1.526cm,draw=black!65] (n3) at (4.340,0.287) {定盗窃罪};
\node[text width=0.500cm,draw=black!65] (n4) at (2.920,1.494) {有};
\node[text width=2.932cm,draw=black!65] (n5) at (4.340,1.494) {第二步：处分意识};
\node[text width=0.500cm,draw=black!65] (n6) at (7.760,1.092) {无};
\node[text width=1.526cm,draw=black!65] (n7) at (9.180,1.092) {定盗窃罪};
\node[text width=0.500cm,draw=black!65] (n8) at (7.760,1.896) {有};
\node[text width=1.526cm,draw=black!65] (n9) at (9.180,1.896) {定诈骗罪};
\draw[edge] (n2.east)--(4.190,0.287)--(4.190,0.287)--(n3.west);
\draw[edge] (n1.east)--(2.770,1.092)--(2.770,0.287)--(n2.west);
\draw[edge] (n6.east)--(9.030,1.092)--(9.030,1.092)--(n7.west);
\draw[edge] (n5.east)--(7.610,1.494)--(7.610,1.092)--(n6.west);
\draw[edge] (n8.east)--(9.030,1.896)--(9.030,1.896)--(n9.west);
\draw[edge] (n5.east)--(7.610,1.494)--(7.610,1.896)--(n8.west);
\draw[edge] (n4.east)--(4.190,1.494)--(4.190,1.494)--(n5.west);
\draw[edge] (n1.east)--(2.770,1.092)--(2.770,1.494)--(n4.west);
\end{tikzpicture}

